\vspace{0.5em}\noindent\textbf{Một chút ``thu hoạch'' sau hành trình làm Software Engineer: từ
viết code đến lúc tự tay làm hệ thống sập rồi tự cứu
lại}\par

\vspace{0.5em}\noindent\textbf{Giới thiệu}\par

Trước đây, tôi từng nghĩ công việc của một Software Engineer chủ yếu là
tiếp nhận yêu cầu, viết code, hoàn thiện tính năng và đưa source code
lên repository. Tuy nhiên, sau khi trực tiếp tham gia phát triển một hệ
thống tương đối đầy đủ, tôi nhận ra rằng viết được code chỉ là bước đầu
tiên.

Phần khó hơn nằm ở việc làm sao để đoạn code đó có thể hoạt động ổn định
trong một hệ thống thật, nơi có nhiều thành phần phụ thuộc lẫn nhau như
frontend, backend API, database, message queue, container, Kubernetes,
CI/CD và hạ tầng AWS.

Trong quá trình phát triển AI-Powered Internship Application Tracker,
tôi đã có cơ hội làm việc với FastAPI, React, PostgreSQL, Redis,
DynamoDB, Amazon SQS, Docker, Kubernetes, GitHub Actions và nhiều dịch
vụ AWS khác. Quá trình này không chỉ giúp tôi học thêm công nghệ mà còn
thay đổi cách tôi nhìn nhận việc xây dựng, triển khai và vận hành phần
mềm.

\vspace{0.5em}\noindent\textbf{Chạy được ở local chưa có nghĩa là chạy được trong
production}\par

Trong môi trường local, các thành phần thường được đặt gần nhau và dễ
kiểm soát. Database có thể chạy trong Docker Compose, Redis nằm trong
cùng network, còn các biến môi trường được khai báo trong file
\texttt{.\allowbreak{}env}.

Tuy nhiên, khi hệ thống được chuyển sang Kubernetes hoặc AWS, cách các
thành phần giao tiếp trở nên phức tạp hơn. Một service muốn kết nối với
service khác cần sử dụng đúng DNS, namespace, port và network
configuration. Ứng dụng muốn kết nối đến database cần có connection
string, route và security group phù hợp. Các pod nằm trong private
subnet cũng cần có đường outbound thích hợp để truy cập Internet hoặc
các dịch vụ bên ngoài.

Từ đó, tôi hiểu rằng khả năng vận hành của một ứng dụng không chỉ phụ
thuộc vào source code mà còn phụ thuộc vào toàn bộ môi trường bao quanh
nó.

\vspace{0.5em}\noindent\textbf{Docker và bài học về môi trường
chạy}\par

Docker giúp chuẩn hóa môi trường chạy giữa các máy, nhưng việc container
hóa hệ thống cũng phát sinh nhiều vấn đề mới.

Một container có thể đang chạy nhưng health check vẫn thất bại. Ứng dụng
có thể khởi động trước khi database sẵn sàng. Biến môi trường có thể tồn
tại trên máy nhưng chưa được truyền đúng vào container. Ngoài ra, các
vấn đề liên quan đến port, volume, Docker Desktop và WSL cũng có thể làm
hệ thống không hoạt động dù code hoàn toàn không thay đổi.

Những trải nghiệm này giúp tôi hình thành thói quen kiểm tra không chỉ
code mà còn cả môi trường thực thi:

\begin{itemize}
\tightlist
\item
  Container đang nhận những biến môi trường nào?
\item
  Dependency đã sẵn sàng chưa?
\item
  Các service có nằm trong cùng network không?
\item
  Port có được expose và mapping chính xác không?
\item
  Health check đang kiểm tra hành vi nào của ứng dụng?
\end{itemize}

\vspace{0.5em}\noindent\textbf{CI/CD không phải là một file YAML thần
kỳ}\par

GitHub Actions mang lại khả năng tự động hóa lint, test, build image,
scan security và deployment. Tuy nhiên, pipeline chỉ hoạt động khi tất
cả giả định bên trong nó đều chính xác.

Workflow có thể thất bại do thiếu secret, sai điều kiện branch, image
tag không đồng nhất hoặc runner thiếu công cụ. Một job cũng có thể bị bỏ
qua vì biến môi trường hoặc điều kiện thực thi không đúng.

Qua quá trình xử lý các workflow thất bại, tôi nhận ra rằng CI/CD không
chỉ là viết một file cấu hình. Người phát triển cần hiểu rõ từng bước
đang làm gì, dữ liệu nào được truyền giữa các job và hệ thống đích cần ở
trạng thái nào trước khi deployment được thực hiện.

\vspace{0.5em}\noindent\textbf{IAM và nguyên tắc không đoán quyền truy
cập}\par

AWS IAM là một trong những phần khiến tôi mất nhiều thời gian nhất. Ban
đầu, tôi nghĩ rằng khi gặp lỗi thiếu quyền thì chỉ cần thêm một policy
có action tương ứng.

Tuy nhiên, quyền truy cập AWS còn phụ thuộc vào trust policy, principal,
condition, repository, branch, resource ARN và cách từng AWS action hỗ
trợ resource-level permission.

Một policy nhìn có vẻ chặt chẽ chưa chắc đã hoạt động đúng. Ví dụ, một
số action bắt buộc phải sử dụng resource \texttt{"*"}, trong khi các
action khác có thể giới hạn vào một ARN cụ thể.

Thay vì tiếp tục phỏng đoán, tôi bắt đầu sử dụng công cụ mô phỏng policy
để kiểm tra từng action. Bài học quan trọng được rút ra là:

\begin{quote}
Không nên đoán quyền truy cập. Cần kiểm chứng chính xác action nào được
phép và action nào đang bị từ chối.
\end{quote}

\vspace{0.5em}\noindent\textbf{Kubernetes: trạng thái Running chưa chắc có nghĩa là hệ
thống
khỏe}\par

Kubernetes có thể hiển thị một pod ở trạng thái \texttt{Running}, nhưng
điều đó chỉ cho biết container vẫn đang tồn tại. Nó không đảm bảo rằng
ứng dụng bên trong đang hoạt động chính xác.

Pod có thể không kết nối được Redis, migration chưa chạy, API không phản
hồi hoặc service không có kết nối outbound. Vì vậy, ngoài trạng thái
pod, cần kiểm tra thêm:

\begin{itemize}
\tightlist
\item
  Application logs.
\item
  Readiness và liveness probes.
\item
  Kubernetes Service và Ingress.
\item
  Database migration.
\item
  Network connectivity.
\item
  Các dependency bên dưới.
\end{itemize}

Điều này giúp tôi hiểu rằng không nên chỉ quan sát trạng thái hạ tầng.
Cần kiểm tra hành vi thật của ứng dụng thông qua endpoint và luồng người
dùng thực tế.

\vspace{0.5em}\noindent\textbf{Sự cố NAT Gateway và bài học về
network}\par

Một trong những sự cố đáng nhớ nhất là khi tôi tắt NAT Gateway để giảm
chi phí. Sau đó, các pod trong private subnet không còn truy cập được
Internet, một số kết nối đến AWS services bị gián đoạn và quá trình
deployment bắt đầu thất bại.

Do code, IAM và security group không thay đổi, nguyên nhân ban đầu khá
khó nhận ra. Sau khi kiểm tra lại network architecture, tôi phát hiện
route outbound đã bị mất vì NAT Gateway không còn hoạt động.

Sự cố này giúp tôi hiểu rõ hơn về public subnet, private subnet,
Internet Gateway, NAT Gateway, route table và VPC Endpoint. Những khái
niệm trước đây chỉ tồn tại trên sơ đồ kiến trúc đã trở thành kiến thức
thực tế sau khi tôi trực tiếp gặp và khắc phục sự cố.

\vspace{0.5em}\noindent\textbf{Database, concurrency và
idempotency}\par

Nhiều lỗi backend không xuất hiện khi chỉ có một người sử dụng hệ thống.
Chúng thường xuất hiện khi hai hoặc nhiều request được gửi gần như đồng
thời.

Nếu hai request cùng kiểm tra một bản ghi chưa tồn tại rồi cùng insert,
dữ liệu trùng lặp có thể được tạo ra. Để hạn chế vấn đề này, hệ thống
cần sử dụng:

\begin{itemize}
\tightlist
\item
  Unique constraints trong database.
\item
  Transaction và rollback.
\item
  HTTP 409 cho các trường hợp conflict.
\item
  Optimistic locking.
\item
  Idempotency key cho các thao tác quan trọng.
\end{itemize}

Ví dụ, thao tác nộp hồ sơ ứng tuyển không nên tạo ra hai application chỉ
vì người dùng bấm nút nhiều lần hoặc client tự động retry request.

\vspace{0.5em}\noindent\textbf{Message queue và transactional
outbox}\par

Một vấn đề khác xuất hiện khi hệ thống cần vừa ghi dữ liệu vào database
vừa gửi message lên queue.

Nếu database commit thành công nhưng việc gửi message thất bại, dữ liệu
đã được tạo nhưng event lại bị mất. Nếu message được gửi trước nhưng
database commit thất bại thì trạng thái giữa hai hệ thống cũng không còn
nhất quán.

Transactional outbox giải quyết vấn đề này bằng cách ghi dữ liệu chính
và event vào cùng một transaction. Một dispatcher riêng sau đó đọc các
event trong outbox và gửi chúng lên Amazon SQS. Nếu việc gửi thất bại,
hệ thống có thể retry. Consumer cũng cần hỗ trợ deduplication để cùng
một event không bị xử lý nhiều lần.

Đây là một ví dụ cho thấy nhiều phần quan trọng của Software Engineering
không trực tiếp xuất hiện trên giao diện người dùng, nhưng lại quyết
định độ tin cậy của toàn bộ hệ thống.

\vspace{0.5em}\noindent\textbf{Cách tiếp cận debug có hệ
thống}\par

Điều quan trọng nhất tôi học được không phải là ghi nhớ thật nhiều câu
lệnh, mà là biết cách chia nhỏ một vấn đề.

Khi có lỗi, tôi bắt đầu kiểm tra theo từng tầng:

\begin{enumerate}
\def\labelenumi{\arabic{enumi}.}
\tightlist
\item
  Lỗi nằm ở application hay infrastructure?
\item
  Lỗi xuất hiện trước hay sau deployment?
\item
  Container và pod có đang chạy không?
\item
  Service có expose đúng port không?
\item
  DNS có resolve được không?
\item
  Network có route được không?
\item
  IAM có cho phép action cần thiết không?
\item
  Database và dependency đã sẵn sàng chưa?
\item
  Bước cuối cùng hoạt động thành công là bước nào?
\end{enumerate}

Thay vì thay đổi nhiều thứ cùng lúc, tôi cố gắng xác định điểm gãy đầu
tiên trong chuỗi xử lý. Cách tiếp cận này giúp quá trình debug có định
hướng và tiết kiệm thời gian hơn.

\vspace{0.5em}\noindent\textbf{Kết luận}\par

Sau hành trình này, cách tôi nhìn nhận Software Engineering đã thay đổi
đáng kể.

Trước đây, khi hệ thống xảy ra lỗi, phản xạ đầu tiên của tôi là sửa
code. Hiện tại, tôi hiểu rằng lỗi có thể nằm ở configuration,
environment variables, container, network, database, permission,
infrastructure hoặc CI/CD.

Tôi cũng bắt đầu quan tâm đến nhiều tình huống hơn ngoài việc tính năng
có chạy được hay không:

\begin{itemize}
\tightlist
\item
  Điều gì xảy ra khi request được gửi hai lần?
\item
  Nếu database commit nhưng queue thất bại thì sao?
\item
  Nếu hai người cùng cập nhật một dữ liệu thì sao?
\item
  Nếu pod restart thì hệ thống phục hồi như thế nào?
\item
  Nếu deployment thất bại thì có thể rollback hay không?
\end{itemize}

Một Software Engineer không phải là người không bao giờ làm hệ thống xảy
ra lỗi. Quan trọng hơn, đó là người có thể hiểu hệ thống rõ hơn sau mỗi
sự cố, khắc phục lỗi có phương pháp và thiết kế để lỗi đó khó lặp lại
hơn.

\vspace{0.5em}\noindent\textbf{Thông tin bài đăng}\par

\begin{itemize}
\tightlist
\item
  \textbf{Chủ đề:} Bài học thực tế từ quá trình phát triển, triển khai
  và vận hành một hệ thống phần mềm.
\item
  \textbf{Ngày đăng:} 28/07/2026.
\item
  \textbf{Nền tảng:} AWS Study Groups.
\item
  \textbf{Trạng thái:} Pending.
\item
  \textbf{Public link:}
  \href{https://www.facebook.com/groups/awsstudygroupfcj/permalink/2227122848052675/\#}{Blog
  3 trên AWS Study Groups}
\end{itemize}
